\chapter*{Appendix T: Modal Logic Foundations}
\addcontentsline{toc}{chapter}{Appendix T: Modal Logic Foundations}

\section{Overview}

This appendix develops the modal logics that govern the verification predicates applied throughout the text. These logics supply the structural principles for reasoning about curvature monotonicity, rewriting correctness, harmonicity preservation, and reconciliation invariants. Two modal systems are defined: the RSVP modality controlling field evolution and the Polyxan modality governing rewriting and patch algebra.

\section{Syntax}

\subsection{Formulas}

Formulas are generated by:
\[
A ::= p \mid \top \mid \bot \mid A \wedge B \mid A \vee B \mid A \to B \mid \Box_{\mathrm{RSVP}} A \mid \Box_{\mathrm{Polyx}} A.
\]

Atomic propositions $p$ represent base invariants such as curvature-boundedness or rewrite stability.

\subsection{Contexts}

A context $\Gamma$ is a finite multiset of formulas. Sequents have the form:
\[
\Gamma \vdash A.
\]

\section{Hilbert Systems}

\subsection{Propositional Core}

The propositional core contains:
\begin{enumerate}
\item all tautologies of intuitionistic propositional logic,
\item modus ponens,
\item necessitation for each modality (restricted to valid formulas).
\end{enumerate}

\subsection{RSVP Modality Axioms}

The RSVP modality satisfies:
\begin{align*}
\mathrm{(R1)} &\quad \Box_{\mathrm{RSVP}}(A \to B) \to (\Box_{\mathrm{RSVP}}A \to \Box_{\mathrm{RSVP}}B), \\
\mathrm{(R2)} &\quad \Box_{\mathrm{RSVP}}A \to A, \\
\mathrm{(R3)} &\quad \Box_{\mathrm{RSVP}}A \to \Box_{\mathrm{RSVP}}\Box_{\mathrm{RSVP}}A.
\end{align*}

These axioms characterize a reflexive and transitive modality corresponding to non-increasing total energy under field evolution.

\subsection{Polyxan Modality Axioms}

The Polyxan modality satisfies:
\begin{align*}
\mathrm{(P1)} &\quad \Box_{\mathrm{Polyx}}(A \to B) \to (\Box_{\mathrm{Polyx}}A \to \Box_{\mathrm{Polyx}}B), \\
\mathrm{(P2)} &\quad \Box_{\mathrm{Polyx}}A \to A, \\
\mathrm{(P3)} &\quad \Box_{\mathrm{Polyx}}A \to \Box_{\mathrm{Polyx}}\Box_{\mathrm{Polyx}}A, \\
\mathrm{(P4)} &\quad \Box_{\mathrm{Polyx}}A \wedge \Box_{\mathrm{Polyx}}B \leftrightarrow \Box_{\mathrm{Polyx}}(A \wedge B).
\end{align*}

The additional distributivity axiom $(\mathrm{P4})$ reflects the confluence and commutativity of the underlying rewriting system.

\section{Sequent Calculi}

\subsection{Structural Rules}

The following rules are admissible:
\begin{align*}
\text{Weakening:} &\quad 
\frac{\Gamma \vdash A}{\Gamma, B \vdash A}, \\
\text{Contraction:} &\quad
\frac{\Gamma, A, A \vdash B}{\Gamma, A \vdash B}, \\
\text{Exchange:} &\quad
\frac{\Gamma, A, B, \Delta \vdash C}{\Gamma, B, A, \Delta \vdash C}.
\end{align*}

\subsection{Modal Rules}

For each modality, introduction and elimination rules are:

\paragraph{RSVP Introduction}
\[
\frac{\Gamma \vdash A \quad \Gamma\ \text{RSVP-invariant}}
{\Gamma \vdash \Box_{\mathrm{RSVP}}A}.
\]

\paragraph{RSVP Elimination}
\[
\frac{\Gamma \vdash \Box_{\mathrm{RSVP}}A}
{\Gamma \vdash A}.
\]

\paragraph{Polyxan Introduction}
\[
\frac{\Gamma \vdash A \quad \Gamma\ \text{Polyxan-invariant}}
{\Gamma \vdash \Box_{\mathrm{Polyx}}A}.
\]

\paragraph{Polyxan Elimination}
\[
\frac{\Gamma \vdash \Box_{\mathrm{Polyx}}A}
{\Gamma \vdash A}.
\]

\section{Kripke Semantics}

\subsection{Frames}

A semantic frame consists of:
\[
W,\qquad 
R_{\mathrm{RSVP}}, \qquad
R_{\mathrm{Polyx}},
\]
where $W$ is a set of worlds; $R_{\mathrm{RSVP}}$ is reflexive and transitive, and $R_{\mathrm{Polyx}}$ is reflexive, transitive, and commutative under intersection.

\subsection{Valuations}

A valuation $V$ assigns atomic formulas to subsets of $W$, respecting curvature- and rewriting-invariance.

\subsection{Interpretation}

The truth relation $\Vdash$ is defined recursively by:
\[
w \Vdash \Box_{\mathrm{RSVP}}A
\quad\text{iff}\quad
\forall w'\,(w R_{\mathrm{RSVP}} w' \Rightarrow w' \Vdash A),
\]
with an analogous clause for $\Box_{\mathrm{Polyx}}$ using $R_{\mathrm{Polyx}}$.

\section{Fixed-Point Operators}

\subsection{Least and Greatest Fixed Points}

Modal fixed points are included using $\mu$ and $\nu$:
\[
\mu X.A(X), \qquad \nu X.A(X),
\]
subject to monotonicity of $A$.

\subsection{Interpretation}

In the RSVP case, least fixed points correspond to energy-minimising invariants, while greatest fixed points correspond to persistent harmonicity.

In the Polyxan case, least fixed points correspond to terminating rewrite sequences, and greatest fixed points to confluence invariants.

\section{Duality Principles}

\subsection{Galois Connections}

The two modalities form a Galois connection on invariant sets:
\[
\Box_{\mathrm{RSVP}}A \leq B
\quad\text{iff}\quad
A \leq \Diamond_{\mathrm{RSVP}}B,
\]
and similarly for the Polyxan operator.

\subsection{Intermodal Commutativity}

Field evolution and rewriting commute up to invariant-preserving equivalence:
\[
\Box_{\mathrm{RSVP}}\Box_{\mathrm{Polyx}}A
\;\leftrightarrow\;
\Box_{\mathrm{Polyx}}\Box_{\mathrm{RSVP}}A,
\]
reflecting the independence of curvature descent and rewriting confluence.

\section{Soundness and Completeness}

\subsection{Soundness}

Every provable sequent $\Gamma \vdash A$ holds in all Kripke frames satisfying the modal axioms.

\subsection{Completeness}

For each modality, the provable formulas coincide with those true in all frames with the appropriate relational properties.

For the combined system, completeness follows from the independence and commutativity of the relations $R_{\mathrm{RSVP}}$ and $R_{\mathrm{Polyx}}$.

\section{Summary}

The modal framework developed here supplies:
\begin{itemize}
\item the logical substrate for field and rewriting invariants,
\item sequent calculi for derivational reasoning,
\item Kripke semantics for curvature and confluence relations,
\item fixed-point constructs for recursive invariants,
\item dualities needed for reconciliation and normalisation proofs.
\end{itemize}

These modal logics provide the verification backbone for the type-theoretic constructions defined in the subsequent appendix.

